\section{A non-abelian example: permutations of three elements}\label{sec:07-groups:04-permutations-example:1}

\texttt{intGroup} is abelian, so \texttt{id\_\allowbreak{}left} never differs from \texttt{id\_\allowbreak{}right}, nor \texttt{inv\_\allowbreak{}left} from \texttt{inv\_\allowbreak{}right} --- in a commutative group these coincide. Does that mean the left/right split in the definition of \texttt{Group} is mere caution, with nothing forcing it? Settling this needs a group where \texttt{op} genuinely fails to commute, built the same way \texttt{intGroup} was, field by field, with both directions checked honestly.

\subsection{The carrier: bijections of a 3-element set}\label{sec:07-groups:04-permutations-example:2}

The symmetric group \(S_3\) is all bijections of a 3-element set, under composition. Represent a permutation of \texttt{Fin 3 := \{0, 1, 2\}} as a bijective function bundled with its inverse and the proofs that they cancel --- an arbitrary \texttt{Fin 3 → Fin 3} need not be a bijection at all.

\begin{lstlisting}
structure Perm3 where
  toFun : Fin 3 (*$\rightarrow$*) Fin 3
  invFun : Fin 3 (*$\rightarrow$*) Fin 3
  left_inv : (*$\forall$*) x, invFun (toFun x) = x
  right_inv : (*$\forall$*) x, toFun (invFun x) = x
\end{lstlisting}

\subsection{The group operation: composition}\label{sec:07-groups:04-permutations-example:3}

\begin{lstlisting}
def Perm3.comp (f g : Perm3) : Perm3 where
  toFun := f.toFun (*$\circ$*) g.toFun
  invFun := g.invFun (*$\circ$*) f.invFun
  left_inv := by
    intro x
    show g.invFun (f.invFun (f.toFun (g.toFun x))) = x
    rw [f.left_inv]
    exact g.left_inv x
  right_inv := by
    intro x
    show f.toFun (g.toFun (g.invFun (f.invFun x))) = x
    rw [g.right_inv]
    exact f.right_inv x
\end{lstlisting}

\texttt{Perm3.comp f g} applies \texttt{g} first, then \texttt{f}, the standard convention. Its inverse composes the two inverses in the opposite order (\texttt{g.invFun ∘ f.invFun}): to undo ``first \(g\), then \(f\),'' first undo \(f\), then undo \(g\). This is \((fg)^{-1} = g^{-1}f^{-1}\), the fact Theorem 3 of Chapter 8 (\texttt{inv\_\allowbreak{}op}) proves abstractly for every group; here it is visible in the construction itself.

\begin{lstlisting}
def Perm3.identity : Perm3 where
  toFun := fun x => x
  invFun := fun x => x
  left_inv := fun _ => rfl
  right_inv := fun _ => rfl
\end{lstlisting}

\begin{lstlisting}
def Perm3.inv (f : Perm3) : Perm3 where
  toFun := f.invFun
  invFun := f.toFun
  left_inv := f.right_inv
  right_inv := f.left_inv
\end{lstlisting}

\texttt{Perm3.inv} swaps \texttt{toFun} with \texttt{invFun}, and correspondingly swaps which proof field plays which role: inverting a bijection swaps which direction counts as forward.

\subsection{Two concrete permutations, and a computed proof they do not commute}\label{sec:07-groups:04-permutations-example:4}

\begin{lstlisting}
-- Swap 0 and 1, leave 2 fixed.
def swap01 : Perm3 where
  toFun := fun x => match x with
    | 0 => 1 | 1 => 0 | 2 => 2
  invFun := fun x => match x with
    | 0 => 1 | 1 => 0 | 2 => 2
  left_inv := by intro x; match x with | 0 => rfl | 1 => rfl | 2 => rfl
  right_inv := by intro x; match x with | 0 => rfl | 1 => rfl | 2 => rfl
\end{lstlisting}

\begin{lstlisting}
-- The 3-cycle 0 (*$\rightarrow$*) 1 (*$\rightarrow$*) 2 (*$\rightarrow$*) 0.
def cycle012 : Perm3 where
  toFun := fun x => match x with
    | 0 => 1 | 1 => 2 | 2 => 0
  invFun := fun x => match x with
    | 0 => 2 | 1 => 0 | 2 => 1
  left_inv := by intro x; match x with | 0 => rfl | 1 => rfl | 2 => rfl
  right_inv := by intro x; match x with | 0 => rfl | 1 => rfl | 2 => rfl
\end{lstlisting}

\begin{lstlisting}
#eval (Perm3.comp swap01 cycle012).toFun 0   -- 0
#eval (Perm3.comp cycle012 swap01).toFun 0   -- 2
\end{lstlisting}

Applying \texttt{cycle012} then \texttt{swap01} sends \(0 \to 1 \to 0\); applying \texttt{swap01} then \texttt{cycle012} sends \(0 \to 1 \to 2\) --- \texttt{\#eval} computes \texttt{0} and \texttt{2}. This is a computed proof, not suggestive evidence, that \texttt{Perm3.comp swap01 cycle012 ≠ Perm3.comp cycle012 swap01}: the group is non-abelian. Chapter 9, Section 7 uses the same compute-a-counterexample move for matrices.

\subsection{\texorpdfstring{Assembling \texttt{Group\ Perm3}}{Assembling Group Perm3}}\label{sec:07-groups:04-permutations-example:5}

\begin{lstlisting}
theorem Perm3.ext {f g : Perm3} (h : (*$\forall$*) x, f.toFun x = g.toFun x)
    (h' : (*$\forall$*) x, f.invFun x = g.invFun x) : f = g := by
  cases f
  cases g
  simp only [mk.injEq]
  constructor
  (*$\cdot$*) funext x; exact h x
  (*$\cdot$*) funext x; exact h' x

def perm3Group : Group Perm3 where
  op := Perm3.comp
  id := Perm3.identity
  inv := Perm3.inv
  assoc := by
    intro f g h
    apply Perm3.ext
    (*$\cdot$*) intro x; rfl
    (*$\cdot$*) intro x; rfl
  id_left := by
    intro f
    apply Perm3.ext
    (*$\cdot$*) intro x; rfl
    (*$\cdot$*) intro x; rfl
  id_right := by
    intro f
    apply Perm3.ext
    (*$\cdot$*) intro x; rfl
    (*$\cdot$*) intro x; rfl
  inv_left := by
    intro f
    apply Perm3.ext
    (*$\cdot$*) intro x; exact f.left_inv x
    (*$\cdot$*) intro x; exact f.left_inv x
  inv_right := by
    intro f
    apply Perm3.ext
    (*$\cdot$*) intro x; exact f.right_inv x
    (*$\cdot$*) intro x; exact f.right_inv x
\end{lstlisting}

\texttt{Perm3.ext} is an extensionality lemma: two \texttt{Perm3}s are equal exactly when their \texttt{toFun}s and \texttt{invFun}s each agree everywhere, the only two fields that matter (proof fields do not, by proof irrelevance, Chapter 6). Each \texttt{Group} field reduces to this check; most cases are \texttt{rfl} directly, and \texttt{inv\_\allowbreak{}left}/\texttt{inv\_\allowbreak{}right} cite exactly the proof obligations already bundled into \texttt{Perm3}.

\begin{mathreading}

\texttt{Perm3} with composition is \(S_3\), the symmetric group on three letters, order \(6\), the smallest non-abelian group. \texttt{swap01} and \texttt{cycle012} are a transposition and a 3-cycle; together they generate all of \(S_3\), as in the standard presentation \(S_3 = \langle r, s \mid r^3 = s^2 = e,\ srs = r^{-1} \rangle\), with \(r\) = \texttt{cycle012} and \(s\) = \texttt{swap01}.

\end{mathreading}

\textbf{Mathlib equivalent.} All of \texttt{Perm3}/\texttt{Perm3.comp}/\texttt{Perm3.ext}/ \texttt{perm3Group} builds one thing: the group of bijections of a 3-element set. \href{https://loogle.lean-lang.org/?q=Equiv.Perm}{\texttt{Equiv.Perm}} is that model, ready-made, a \href{https://loogle.lean-lang.org/?q=Group}{\texttt{Group}} instance for any type.

\begin{lstlisting}
example : Group (Equiv.Perm (Fin 3)) := inferInstance

-- Swap 0 and 1, leave 2 fixed (*\textemdash{}*) the Mathlib analogue of `swap01`.
def swap01' : Equiv.Perm (Fin 3) := Equiv.swap 0 1

-- The 3-cycle 0 (*$\rightarrow$*) 1 (*$\rightarrow$*) 2 (*$\rightarrow$*) 0 (*\textemdash{}*) the Mathlib analogue of `cycle012`.
def cycle012' : Equiv.Perm (Fin 3) := finRotate 3

#eval (swap01' * cycle012') 0   -- 0
#eval (cycle012' * swap01') 0    -- 2
\end{lstlisting}

No \texttt{Perm3} bundle, no hand-written extensionality lemma. \texttt{Equiv.Perm (Fin 3)}, the type of bijections \texttt{Fin 3 ≃ Fin 3}, is already a group; \href{https://loogle.lean-lang.org/?q=Equiv.swap}{\texttt{Equiv.swap}} and \href{https://loogle.lean-lang.org/?q=finRotate}{\texttt{finRotate}} construct a transposition and a rotation, \texttt{*} is the registered operation matching \texttt{Perm3.comp}'s convention. The two \texttt{\#eval}s are the same non-commutativity witness, now on the library's \(S_3\).

